\documentclass[11pt,a4paper]{article}
\usepackage[margin=0.6in]{geometry}
\usepackage{enumitem}
\usepackage{titlesec}
\usepackage{hyperref}
\usepackage{xcolor}
\usepackage{parskip}
\pagestyle{empty}
\setlength{\parindent}{0pt}
\definecolor{headingcolor}{RGB}{0,51,102}
\hypersetup{colorlinks=true,urlcolor=headingcolor}
\titleformat{\section}{\large\bfseries\color{headingcolor}}{}{0em}{}[\titlerule]
\titlespacing*{\section}{0pt}{10pt}{4pt}
\begin{document}
\begin{center}
{\Huge\bfseries Maya Chen}\\[6pt]
{\small maya.chen@example.com \quad | \quad +1 202-555-0147}
\end{center}
\section{Education}
\textbf{Northbridge University (Fictional)} \hfill \textit{Sep. 2022 - Jun. 2026}\\
BSc in Computer Science \hfill Boston, MA\\
\textit{GPA: 3.91/4.00, Class Ranking: Top 5\%}

\textbf{Relevant Coursework:} Machine Learning, Algorithms, Probability, Linear Algebra
\vspace{4pt}

\section{Research Experience}
\textbf{Undergraduate Research Assistant - Robust Learning} \hfill \textit{Sep. 2024 - Jun. 2026}\\
\textit{Northbridge Reliable AI Lab | Advisor: Dr. Elena Brooks}
\begin{itemize}[leftmargin=*,noitemsep,topsep=2pt]
\item Implemented image-classification baselines in PyTorch and evaluated robustness under distribution shifts.
\item Compared augmentation and calibration methods across 12 corruption settings using five random seeds.
\item Reduced expected calibration error from 8.4\% to 5.1\% on the held-out evaluation split.
\item Defined corruption severity groups and kept preprocessing, checkpoints and inference settings fixed across comparisons.
\item Conducted ablations on label smoothing, augmentation strength and temperature scaling to separate their effects.
\item Reported clean-data accuracy alongside robustness metrics to make accuracy-calibration tradeoffs explicit.
\item Wrote data-loading tests and saved configurations, environment versions and seed-level results for replication.
\end{itemize}

\textbf{Summer Research Student - Dataset Quality} \hfill \textit{Jun. 2024 - Aug. 2024}\\
\textit{Northbridge Data Science Center | Advisor: Dr. Daniel Reed}
\begin{itemize}[leftmargin=*,noitemsep,topsep=2pt]
\item Audited 18,000 labels and documented annotation disagreements in an internal research report.
\item Created stratified data splits and leakage checks for reproducible model comparison.
\item Designed an annotation review protocol covering ambiguous labels, near-duplicates and class imbalance.
\item Reviewed disagreements with two student annotators and recorded decisions without overwriting the original labels.
\item Compared class distributions before and after filtering and identified categories with insufficient coverage.
\end{itemize}

\section{Professional Experience}
\textbf{Cedar Analytics (Fictional)} \hfill \textit{Jun. 2025 - Aug. 2025}\\
\textit{Machine Learning Intern} \hfill Remote
\begin{itemize}[leftmargin=*,noitemsep,topsep=2pt]
\item Built feature-validation checks and experiment summaries for an internal demand-forecasting prototype.
\item Implemented time-based validation and compared seasonal baselines before testing tree-based forecasting models.
\item Investigated missing values, delayed records and feature drift, documenting risks for deployment.
\item Prepared a handover notebook explaining evaluation assumptions and limitations to the engineering team.
\end{itemize}
\vspace{4pt}

\section{Selected Projects}
\textbf{\href{https://example.com/demos/reproducibility}{Experiment Reproducibility Toolkit}} \hfill \textit{2025}\\
\textit{Python, PyTorch, YAML, Git}
\begin{itemize}[leftmargin=*,noitemsep,topsep=2pt]
\item Built configuration-based runners, seed tracking and aggregated confidence-interval reports.
\item Documented benchmark setup and added tests for split integrity and missing metrics.
\item Added resumable execution and explicit failure records to avoid silently excluding unsuccessful experiments.
\item Created a tutorial with a small synthetic dataset and a worked example of interpreting uncertainty.
\item Separated dataset configuration from training code and checked repeated runs for consistent split membership.
\end{itemize}

\textbf{\href{https://example.com/demos/calibration-dashboard}{Calibration Diagnostics Dashboard}} \hfill \textit{Jan. 2026 - May 2026}\\
\textit{Python, Streamlit, pandas, Plotly}
\begin{itemize}[leftmargin=*,noitemsep,topsep=2pt]
\item Built reliability diagrams and class-level error summaries from saved prediction files.
\item Added subgroup filters and displayed sample sizes to discourage conclusions from sparse categories.
\item Compared expected calibration error under multiple binning schemes and documented estimator sensitivity.
\item Wrote tests for malformed probability vectors, missing labels and classes absent from the evaluation split.
\end{itemize}

\textbf{\href{https://example.com/demos/small-data-learning-study}{Small-Data Learning Study}} \hfill \textit{2024}\\
\textit{Python, PyTorch, scikit-learn}
\begin{itemize}[leftmargin=*,noitemsep,topsep=2pt]
\item Designed a course study comparing linear models, shallow networks and pretrained representations under restricted label budgets.
\item Evaluated five training-set sizes with fixed class balance and repeated subsampling.
\item Compared validation-based model selection with a deliberately overfit baseline to demonstrate evaluation leakage.
\item Recorded compute time and parameter count alongside classification accuracy.
\item Wrote a project report explaining why gains on one dataset did not establish general superiority.
\end{itemize}

\textbf{\href{https://example.com/demos/model-error-review-workspace}{Model Error Review Workspace}} \hfill \textit{2025}\\
\textit{Python, SQLite, pandas}
\begin{itemize}[leftmargin=*,noitemsep,topsep=2pt]
\item Built a local workspace for reviewing misclassified examples and grouping errors by category.
\item Designed an annotation table retaining model version, confidence, review decision and reviewer notes.
\item Added export checks to prevent duplicate examples from entering both training and evaluation collections.
\item Analyzed error clusters with the research team and proposed targeted data-review tasks.
\item Documented that qualitative error categories were exploratory and required independent validation.
\end{itemize}

\section{Technical Skills}
\textbf{Programming Languages:} Python, C++, SQL

\textbf{Frameworks \& Libraries:} PyTorch, NumPy, pandas, scikit-learn

\textbf{Systems \& Tools:} Linux, Git, Docker

\textbf{Specializations:} Robust learning, experiment design, model calibration

\section{Honors \& Awards}
\textbf{Northbridge Undergraduate Research Scholarship} \hfill 2025 \textit{-- Supported a year-long independent research project.}

\textbf{Faculty Academic Excellence Award} \hfill 2024 \textit{-- Recognized performance in mathematics and computing.}

\textbf{Research Poster Presentation Commendation} \hfill 2026 \textit{-- Recognized a clear presentation of methods and experimental limitations.}
\end{document}
